% =====================================================================
%  zestaw przykładów:
%   - wzór \label{eq:indeksliczenie} z odsyłaczem "(str. X)" w tekście,
%   - listing w Rust (minted),
%   - duży rysunek rysowany w TikZ-u (komórki tensorów),
%   - algorytm w pseudokodzie (algorithm2e).
%  Każdy element ma \label, więc można się do niego odwoływać \ref/\pageref.
% =====================================================================
\section{Implementacja podstawowych komponentów} \label{sect:impl}
Tutaj coś piszemy
\subsection{Mechanizm Strides i adresowanie pamięci}
Mapowanie indeksów, wykorzystywane w wielu operacjach, operuje na mechanizmie
kroków (\textit{strides}). Umożliwia to mapowanie wielowymiarowego indeksu, na przykład $(i, j, k)$
na fizyczny adres -- przesunięcie (\textit{offset}) względem początku wektora danych. Obliczenie przesunięcia zostało
ukazany z użyciem wzoru \ref{eq:indeksliczenie} (str. \pageref{eq:indeksliczenie}).
\begin{equation} \label{eq:indeksliczenie}
	\text{offset} = \sum_{d=0}^{D-1} i_d \cdot s_d
\end{equation}
gdzie $i_d$ to indeks w wymiarze $d$, a $s_d$ to wartość kroku (\textit{stride}) dla tego wymiaru. Operacja ta
jest implementowana w metodzie \texttt{physical\_offset}.

\begin{minted}[
    frame=lines,        % Ramka na górze i na dole
    framesep=2mm,       % Odstęp ramki od kodu
    baselinestretch=1.2,% Interlinia dla czytelności
    fontsize=\small,    % Nieco mniejsza czcionka
    linenos             % Numery linii
]
{rust}
pub fn physical_offset(&self, logical_idx: impl AsRef<[usize]>) -> isize {
  let logical_idx = logical_idx.as_ref();
  logical_idx
    .iter()
    .zip(self.strides.iter())
    .fold(0isize, |acc, (&ind, &stride)| acc + ind as isize * stride)
}
\end{minted}


\begin{figure}[h]
	\centering
	\begin{tikzpicture}[every node/.style={font=\scriptsize}]

		% DEFINICJE KOLORÓW
		\definecolor{custorange}{RGB}{255,179,102}
		\definecolor{custpink}{RGB}{255,153,204}
		\definecolor{custcyan}{RGB}{153,255,255}
		\definecolor{custlime}{RGB}{204,255,153}

		% PARAMETRY
		\pgfmathsetmacro{\cellsize}{1.0}
		\pgfmathsetmacro{\depthshift}{1.2}
		\pgfmathsetmacro{\imageshift}{1}
		\pgfmathsetmacro{\imagespace}{6.5}
		\pgfmathsetmacro{\matrixstart}{-6.0}

		% 1. zdjęcia

		\begin{scope}[shift={(\imageshift, 0)}]
			\foreach \h in {0,1} {
					\foreach \w in {0,1,2} {
							\pgfmathsetmacro{\xpos}{\w*\cellsize + \depthshift}
							\pgfmathsetmacro{\ypos}{-\h*\cellsize - \depthshift}

							\fill[custpink, opacity=0.8] (\xpos, \ypos) rectangle ++(\cellsize, -\cellsize);
							\draw[black] (\xpos, \ypos) rectangle ++(\cellsize, -\cellsize);
							\node[font=\tiny] at (\xpos + 0.5*\cellsize, \ypos - 0.5*\cellsize)
							{$i_{1,\w,\h,0}$};
						}
				}

			\foreach \h in {0,1} {
					\foreach \w in {0,1,2} {
							\pgfmathsetmacro{\xpos}{\w*\cellsize}
							\pgfmathsetmacro{\ypos}{-\h*\cellsize}

							\fill[custorange, opacity=0.8] (\xpos, \ypos) rectangle ++(\cellsize, -\cellsize);
							\draw[black] (\xpos, \ypos) rectangle ++(\cellsize, -\cellsize);
							\node[font=\tiny] at (\xpos + 0.5*\cellsize, \ypos - 0.5*\cellsize)
							{$i_{0,\w,\h,0}$};
						}
				}

			\draw[dashed, gray] (0,0) -- (\depthshift, -\depthshift);
			\draw[dashed, gray] (3*\cellsize,0) -- (3*\cellsize+\depthshift, -\depthshift);
			\draw[dashed, gray] (0,-2*\cellsize) -- (\depthshift, -2*\cellsize-\depthshift);
			\draw[dashed, gray] (3*\cellsize,-2*\cellsize) -- (3*\cellsize+\depthshift, -2*\cellsize-\depthshift);

			\node[above, font=\small\bfseries] at (\cellsize + 0.5, 0.3) {Zdjęcie 0 ($B=0$)};
		\end{scope}

		\begin{scope}[shift={(\imagespace, 0)}]
			\foreach \h in {0,1} {
					\foreach \w in {0,1,2} {
							\pgfmathsetmacro{\xpos}{\w*\cellsize + \depthshift}
							\pgfmathsetmacro{\ypos}{-\h*\cellsize - \depthshift}

							\fill[custlime, opacity=0.8] (\xpos, \ypos) rectangle ++(\cellsize, -\cellsize);
							\draw[black] (\xpos, \ypos) rectangle ++(\cellsize, -\cellsize);
							\node[font=\tiny] at (\xpos + 0.5*\cellsize, \ypos - 0.5*\cellsize)
							{$i_{1,\w,\h,1}$};
						}
				}

			\foreach \h in {0,1} {
					\foreach \w in {0,1,2} {
							\pgfmathsetmacro{\xpos}{\w*\cellsize}
							\pgfmathsetmacro{\ypos}{-\h*\cellsize}

							\fill[custcyan, opacity=0.8] (\xpos, \ypos) rectangle ++(\cellsize, -\cellsize);
							\draw[black] (\xpos, \ypos) rectangle ++(\cellsize, -\cellsize);
							\node[font=\tiny] at (\xpos + 0.5*\cellsize, \ypos - 0.5*\cellsize)
							{$i_{0,\w,\h,1}$};
						}
				}

			\draw[dashed, gray] (0,0) -- (\depthshift, -\depthshift);
			\draw[dashed, gray] (3*\cellsize,0) -- (3*\cellsize+\depthshift, -\depthshift);
			\draw[dashed, gray] (0,-2*\cellsize) -- (\depthshift, -2*\cellsize-\depthshift);
			\draw[dashed, gray] (3*\cellsize,-2*\cellsize) -- (3*\cellsize+\depthshift, -2*\cellsize-\depthshift);

			\node[above, font=\small\bfseries] at (\cellsize + 0.5, 0.3) {Zdjęcie 1 ($B=1$)};
		\end{scope}

		\draw[->, ultra thick] (6, -3.5) -- (6, \matrixstart + 0.5)
		node[midway, right] {\texttt{reshape}};

		% 2. MACIERZ

		\begin{scope}[shift={(0, \matrixstart)}]
			\foreach \x/\cellcolor/\idx in {
			0/custorange/i_{0,0,0,0}, 1/custorange/i_{0,1,0,0},
			2/custorange/i_{0,2,0,0}, 3/custorange/i_{0,0,1,0},
			4/custorange/i_{0,1,1,0}, 5/custorange/i_{0,2,1,0},
			6/custcyan/i_{0,0,0,1}, 7/custcyan/i_{0,1,0,1},
			8/custcyan/i_{0,2,0,1}, 9/custcyan/i_{0,0,1,1},
			10/custcyan/i_{0,1,1,1}, 11/custcyan/i_{0,2,1,1}
			} {
			\fill[\cellcolor, opacity=0.8] (\x*\cellsize, 0) rectangle ++(\cellsize, -\cellsize);
			\draw[black] (\x*\cellsize, 0) rectangle ++(\cellsize, -\cellsize);
			\node[font=\tiny] at (\x*\cellsize + 0.5*\cellsize, -0.5*\cellsize) {$\idx$};
			}
		\end{scope}

		\begin{scope}[shift={(0, \matrixstart - 1*\cellsize)}]
			\foreach \x/\cellcolor/\idx in {
			0/custpink/i_{1,0,0,0}, 1/custpink/i_{1,1,0,0},
			2/custpink/i_{1,2,0,0}, 3/custpink/i_{1,0,1,0},
			4/custpink/i_{1,1,1,0}, 5/custpink/i_{1,2,1,0},
			6/custlime/i_{1,0,0,1}, 7/custlime/i_{1,1,0,1},
			8/custlime/i_{1,2,0,1}, 9/custlime/i_{1,0,1,1},
			10/custlime/i_{1,1,1,1}, 11/custlime/i_{1,2,1,1}
			} {
			\fill[\cellcolor, opacity=0.8] (\x*\cellsize, 0) rectangle ++(\cellsize, -\cellsize);
			\draw[black] (\x*\cellsize, 0) rectangle ++(\cellsize, -\cellsize);
			\node[font=\tiny] at (\x*\cellsize + 0.5*\cellsize, -0.5*\cellsize) {$\idx$};
			}
		\end{scope}


	\end{tikzpicture}
	\caption{Przedstawienie zmiany kształtu tensora zbioru dwóch zdjęć na
		reprezentację macierzową bez modyfikacji danych właściwych, a jedynie metadanych.
		Źródło: opracowanie własne}
	\label{fig:conv2dimgreshape}
\end{figure}

\begin{algorithm}[h]
	% UWAGA na kolejność: najpierw \caption, potem \label — inaczej
	% odsyłacz \ref{alg:tensor-iter} pokaże złą numerację.
	\caption{Iterator po tensorze (column-major order)}
	\label{alg:tensor-iter}
	\KwIn{Struktura iteratora $I$}
	\KwOut{Następna wartość lub $\bot$ (brak elementów)}

	\tcp{Struktura TensorIter:}
	\tcp{$I.\text{view}$ -- referencja do tensora}
	\tcp{$I.\text{logical\_pos}$ -- bieżąca pozycja [i₀, i₁, ...]}
	\tcp{$I.\text{finished}$ -- flaga zakończenia}

	\BlankLine
	\If{$I.\text{finished}$}{
		\Return $\bot$
	}

	\tcp{Oblicz fizyczny offset i pobierz wartość}
	$\text{offset} \gets \sum_{d} I.\text{logical\_pos}[d] \cdot I.\text{view}.\text{strides}[d]$\;
	$\text{value} \gets I.\text{view}.\text{data}[\text{offset}]$\;

	\tcp{Inkrementuj pozycję (column-major)}
	$k \gets 0$\;
	\While{$k < |I.\text{view}.\text{shape}|$}{
		$I.\text{logical\_pos}[k] \gets I.\text{logical\_pos}[k] + 1$\;
		\If{$I.\text{logical\_pos}[k] < I.\text{view}.\text{shape}[k]$}{
			\Return $\text{value}$
		}
		$I.\text{logical\_pos}[k] \gets 0$\;
		$k \gets k + 1$\;
	}

	$I.\text{finished} \gets \text{true}$\;
	\Return $\text{value}$
\end{algorithm}

